\documentclass[10pt,a4paper]{article}
\usepackage[utf8]{inputenc}
\usepackage[a4paper,left=3cm,right=2cm]{geometry}
\usepackage{amsmath}
\usepackage[thmmarks, amsmath, thref]{ntheorem}
\usepackage{amsfonts}
\usepackage{amssymb}
\usepackage{fancyhdr}
\pagestyle{fancy}
\usepackage{anysize}
\usepackage{graphicx}
\usepackage{float}
\usepackage{hyperref}
\usepackage{enumitem}
\usepackage{verbatim}
\newtheorem{theorem}{Theorem}
\theoremsymbol{\ensuremath{\blacksquare}}
\newtheorem*{solution}{Solución:}

\usepackage{xcolor}

% Margenes y formalidades

\textwidth 6.25in % Margen Derecho y ancho del texto
\textheight 8.5in
\oddsidemargin 0in % Margen Izquierdo
\headheight 0in % Largo del texto
\leftmargin 1in
\parindent 0em % Salto entre la indentación (Sangría)https://www.overleaf.com/4519127669sdpptkhvsbfn
\topmargin -0.5in % Margen Superior
\parskip 0.5ex % Salto entre parrafos
\baselineskip 1pt



\lhead{
\includegraphics[scale=0.08]{Figuras/logoUSM.png}} % texto izquierda de la cabecera
%\chead{}                                      % texto centro de la cabecera
\rhead{Universidad Técnica Federico Santa María\\Ingeniería Comercial } % número de página a la derecha
\renewcommand{\headrulewidth}{0.4pt} % grosor de la línea de la cabecera
\renewcommand{\footrulewidth}{0.4pt} % grosor de la línea del pie

\begin{document}
\vspace{5cm}



\begin{center}
\begin{large}
\textbf{\\CERTAMEN 2}
\end{large}\\
\end{center}

\begin{center}
\begin{large}
\textbf{Econometría ICS-294}
\end{large}\\
\vspace{0.2cm}
I Semestre 2025\\
Profesores: Jocelyn Tapia Stefanoni, Franco Antonucci Fernández y Ricardo Pruneda Urz\'ua\\
%Ayudante: 
\end{center}

\textbf{Instrucciones:} 

\begin{itemize}

    \item Usted dispone de 90 minutos  para intentar los 100 puntos del certamen.
    \item Cuide la presentación y redacción de sus preguntas.
    \item No puede utilizar celular, computador o tablet, pero sí una calculadora TI.
     
\end{itemize}

\section{RdA3.2  Preguntas cortas (40 puntos)}
{\bf Elija solo 3 preguntas y responda. }
\begin{enumerate}
   \item    Suponga un modelo dado por:
    
    $$\ln{wage} = \beta_0 + \beta_1 educ + \beta_2 exper + \mu$$
    
    Donde $educ$: educación; $exper$: años de experiencia. Al correr una regresión del modelo se obtienen los siguientes valores
    \begin{table}[H] 
    \centering
    \caption{Resultados de la regresión}
    \begin{tabular}{lcc}
    \hline
    \hline
    \textbf{Variable} & \textbf{Coeficiente} & \textbf{Error Estándar} \\
    \hline
    Intercepto & 1.20$^{*}$ & 0.15 \\
    educ & 0.08$^{***}$ & 0.01 \\
    exper & 0.03$^{***}$ & 0.005 \\
    \hline
    \multicolumn{3}{l}{\textbf{n} = 1000} \\
    \multicolumn{3}{l}{$R^2 = 0.32$} \\
    \multicolumn{3}{l}{$F = 235.4$ (gl = 2,997)} \\
    \hline
    \hline
    \textit{Nota:} $^{*}p<0.1$; $^{**}p<0.05$; $^{***}p<0.01$
    \end{tabular}
    \end{table}
    
    {Con lo cual un investigador le dice a usted que no hay retornos positivos por más experiencia al 5\% de significancia. Comente.}\\
    
    {\color{blue} Lo que dice el investigados es falso. Para poder verificar lo que dice el investigador se puede hacer la siguiente prueba de hipótesis:
    
    $$H_0: \beta_2 = 0$$
    $$H_1: \beta_2 > 0$$
    
    Luego, se calcula el estadístico $t$ como $t_{\hat{\beta_2}} = \frac{0.03}{0.005} = 6$. Se compara ese valor contra el valor crítico. En este caso como hay $n-k-1 = 997$ gl, se pueden usar valores de la normal. El valor crítivo a un nivel de significancia del 5\% es 1.645.
    
    Por lo tanto, como $6>1,645$ se concluye que $\beta_2$ es estadísticamente mayor a 0 a un nivel de significancia del 5\%, es decir, la conclusión del investigador es falsa.}\\


    \item  ¿Cómo podría reducir la varianza de los estimadores MCO en un estudio econométrico real? Señale al menos dos formas.

{\color{blue}
\begin{itemize}
    \item Aumentando el tamaño de muestra $n$, ya que mejora la precisión.
    \item Aumentando la variabilidad de las variables independientes, evitar la colinealidad.
    \item Mejorando la especificación del modelo para reducir la varianza del error ($\sigma^2$).

\end{itemize}}
    \item 
     Un estimador es consistente cuando $n \longrightarrow \infty$; en este caso la distribución del estimador colapsa hasta el valor del parámetro que se está estimando. Esto significa que el parámetro es consistente e insesgado. Comente.\\

    {\color{blue} Falso. La primera parte del enunciado es correcta y, en efecto, cuando eso ocurre el parámetro es consistent. Pero la segunda parte no, no necesariamente es el estimado será insesgado, pues esto último depende de que se cumplan los supuestos de Gauss-Markov.}

    \item   En un modelo de regresión lineal simple estimado usando una base de datos con  150 observaciones, la suma de residuos al cuadrado es 19.000, determine  $\hat{\sigma}^2$.

{\color{blue}
\[
\hat{\sigma}^2 = \frac{SSR}{n - k - 1} = \frac{19000}{150 - 1 - 1} = \frac{19000}{148} \approx 128,38
\]}

\end{enumerate}


\section{RdA3.2  Ejercicios (60 puntos)}


{\bf A continuación elija solo 3 preguntas y responda. }

\begin{enumerate}
 \item  Una investigación se encuentra analizando los principales determinantes del peso de los recién nacidos.
Se dispone de una base de datos con las siguientes variables: 

\begin{itemize}
    \item 
bwght: peso al nacer, en libras.
\item cigs: número promedio de cigarros diarios que fumó la madre durante el embarazo.
\item parity: orden de nacimiento del niño.
\item faminc: ingreso anual de la familia.
\item mo\_educ: años de escolaridad de la madre.
\item fa\_educ: años de escolaridad del padre.
\item male: variable dicotómica que toma el valor 1 si el recién nacido es varón.
\end{itemize}

Los resultados obtenidos están representados en la siguiente ecuación, donde en paréntesis aparecen los errores estándar de los estimadores.
    
{\scriptsize
\begin{equation*}
\begin{array}{ll}
\widehat{log(bwght)}  = & 1,1 -0,0052 cigs+0,011 log( faminc) +0,017 parity+0,034 male-0,0065mo\_educ +0,005fa\_educ\\
& (0,38) \quad (0, 001) \quad \quad \quad \quad (0,004)\quad \quad \quad \quad \quad (0,006) \quad \quad \quad \quad (0,011)\quad \quad \quad \quad  (0,03)\quad \quad \quad \quad \quad (0,0046)\\
 & n =1.191 \quad R^2=  0,493
\end{array}
\end{equation*}}


\begin{enumerate}
    
    \item (5 puntos) Analice el efecto de un aumento del 10\% en el ingreso anual de la familia sobre el  peso al nacer. Interprete este efecto en términos económicos y evalúe su significancia estadística utilizando un nivel de significancia del 1\% (\( \alpha = 0{,}01 \)). ¿Se puede concluir que el ingreso familiar tiene un efecto significativo sobre el peso al nacer?

{\color{blue} El coeficiente de \(\log(\text{faminc})\) es 0{,}011, por lo tanto, un aumento del 10\% en el ingreso familiar incrementa el peso al nacer en:

\[
0{,}011 \times 10 = 0{,}11\%
\]

Este término puede interpretarse como una elasticidad al estar ambas variables medidas en logaritmos.

Para evaluar la significancia estadística:

\[
t = \frac{0{,}011}{0{,}004} = 2{,}75
\]

Como tenemos gl=1191-7= 1184 usamos el valor cr\'itico de la normal. El valor del estad\'istico supera el valor crítico de 2{,}576  para \(\alpha = 0{,}01\), por lo que se concluye que el ingreso familiar tiene un efecto estadísticamente significativo sobre el peso al nacer al 1\%.
}

    \item (5 puntos) Evalúe la significancia estadística individual de las variables \textit{escolaridad del padre} (\( \text{fatheduc} \)) y \textit{escolaridad de la madre} (\( \text{motheduc} \)) en el modelo estimado. Para ello, realice las pruebas t correspondientes y determine si alguna de estas variables tiene un efecto significativo sobre el peso al nacer, utilizando un nivel de significancia del 5\%.

    {\color{blue}Para \textbf{mo\_educ} (madre):

\[
t = \frac{-0{,}0065}{0{,}03} \approx -0{,}217
\]

Para \textbf{fa\_educ} (padre):

\[
t = \frac{0{,}005}{0{,}0046} \approx 1{,}087
\]

Ambos valores son menores en valor absoluto a 1,96 (valor crítico de la normal ).
Por lo tanto, no se rechaza la hipótesis nula para ninguna de las dos variables a un nivel de significancia del 5\%. Ninguna tiene efecto estadísticamente significativo de manera individual.
}

    \item {(10 puntos)} Se presenta a continuación un modelo alternativo estimado para explicar el logaritmo del peso al nacer:

    \begin{equation*}
    \begin{aligned}
    \widehat{\log(\text{bwght})} =\ & 1{,}66 - 0{,}0044\, \text{cigs} + 0{,}0093\, \log(\text{faminc}) + 0{,}016\, \text{parity} + 0{,}027\, \text{male} \\
    & (0{,}22) \quad (0{,}0009) \quad \quad (0{,}0059) \quad \quad (0{,}006) \quad \quad (0{,}010) \\
    & n = 1.191 \quad \quad R^2 = 0{,}472
    \end{aligned}
    \end{equation*}

    Utilizando información de este modelo y del modelo original que incluía las variables \textit{motheduc} y \textit{fatheduc}, realice una prueba de significancia conjunta (prueba F) para evaluar si ambas variables pueden ser excluidas del modelo sin pérdida significativa de capacidad explicativa. ¿Debe la investigación preferir el modelo completo o el modelo reducido? Justifique su elección en función de criterios estadísticos y del objetivo del análisis.

{\color{blue} Modelo no restringido: \(R^2_{ur} = 0{,}493\)

Modelo restingido: \(R^2_r = 0{,}472\)

Prueba F:

\[
F = \frac{(0{,}493 - 0{,}472)/2}{(1 - 0{,}493)/(1191 - 7)} = \frac{0{,}0105}{0{,}000428} \approx 24{,}53
\]

El valor crítico de \(F_{2, 1184}\) al 1\% es aproximadamente 4{,}60. Como \(24{,}53 > 4{,}60\), se rechaza la hipótesis nula.

Aunque las variables no son significativas individualmente, \textbf{conjuntamente sí lo son}. Por tanto, desde un punto de vista estadístico, se prefiere el \textbf{modelo completo}, ya que ofrece una mejor capacidad explicativa.
 }
\end{enumerate}
    \item Considere la siguiente información:
 \[
n=600 \quad \hat{\beta} =  \begin{bmatrix}
1,5 \\
2 \\
3,55 
\end{bmatrix} \quad \hat{\sigma}^2 = 141925; \quad 
(X^\top X)^{-1} = 
\begin{bmatrix}
0.05987 & -0.00177 & -0.00035 \\
-0.00177 & 0.00030 & -0.00002 \\
-0.00035 & -0.00002 & 0.00001
\end{bmatrix}
\]

\begin{enumerate}
    \item (9 puntos) Determine los errores estándar de $\hat{\beta}_0$, $\hat{\beta}_1$ y $\hat{\beta}_2$.
    
  {\color{blue}  
    \[
    \widehat{\text{Var}}(\hat{\beta}_0) = 141925 \times 0.05987 = 8497.04975\\ \]
    
El error estándar: $\hat{\beta}_0$.

 $\sqrt{8497.04975} \approx 92.18$

   \[
    \widehat{\text{Var}}(\hat{\beta}_1) = 141925 \times 0.00030 = 42.5775\\
    \]
    
El error estándar: $\hat{\beta}_1$.

 $\sqrt{42.5775} \approx 6.53$

  \[
    \widehat{\text{Var}}(\hat{\beta}_2) = 141925 \times 0.00001 = 1.41925\\ \]
    
El error estándar: $\hat{\beta}_2$.

 $\sqrt{1.41925} \approx 1.19$}
 
\item(11 puntos) Plantee la hip\'otesis  que testea $\beta_2$ es el doble de  $\beta_1$. 
¿Se puede rechazar la hipótesis nula con la información que disponemos con una significancia del 1\%?
\end{enumerate}

{\color{blue}  \textbf{Formulación de la hipótesis:}

\[
H_0: \beta_2 - 2\beta_1 = 0 \qquad H_1: \beta_2 - 2\beta_1 \neq 0
\]

 \textbf{Construcción del estadístico t:}

\[
t_e = \frac{\hat{\beta}_2 - 2\hat{\beta}_1}{\sqrt{\widehat{\text{Var}}(\hat{\beta}_2 - 2\hat{\beta}_1)}}
\]
 \textbf{Cálculo de la varianza del estimador lineal:}

\[
\widehat{\text{Var}}(\hat{\beta}_2 - 2\hat{\beta}_1) =
\widehat{\text{Var}}(\hat{\beta}_2) + 4 \cdot \widehat{\text{Var}}(\hat{\beta}_1) - 4 \cdot \widehat{\text{Cov}}(\hat{\beta}_2, \hat{\beta}_1)
\]

\[
= 1.41925 + 4 \cdot 42.5775 - 4 \cdot (141925*-0.00002) = 1.41925 + 170.31 + 4\times 2.84 = 183.083
\]

 \textbf{Error estándar:}
\[
se(\hat{\beta}_2 - 2\hat{\beta}_1) = \sqrt{183.083} \approx 13.53
\]

 \textbf{Valor del estadístico t:}
\[
t_e = \frac{3.55 - 2 \cdot 2}{13.53} = \frac{-0.45}{13.53} \approx -0.033
\]



Valor crítico bilateral con gl = 600-2-1: $t_{0.005} \approx \pm 2.576$ (usando el valor de la normal)

Como $|t_e| = 0.22 < 2.576$, no se rechaza $H_0$.

Con un nivel de significancia del 1\%, no se rechaza la hipótesis nula de que $\beta_2 = 2\beta_1$. No hay suficiente evidencia para decir que los efectos son distintos. }

    
    \item     Suponga que, debido al tornado que hubo en Puerto Varas (Región de Los Lagos, Chile) el día Domingo 25 de Mayo de 2025, le piden a usted estimar un modelo de regresión lineal simple de la fuerza de un tornado (medido en la escala Fujita). Usted toma muchos datos de tornados y estima que el modelo es:

    $$F_i = \gamma_0 + \gamma_1 \cdot V_i + \gamma_2 \cdot P_i + \gamma_3 \cdot T_i + u_i$$

    Donde $F$ es la fuerza del tornado; $V$ es la velocidad del viento ($km/h$); $P$: nivel de precipitaciones al momento de ocurrir el tornado (en mm); $T$: temperatura al momento de ocurrir el tornado (en ºC).\\

    Se tiene lo siguiente:

\begin{table}[H] \centering 
  %\caption{Regresión múltiple sobre la fuerza del tornado} 
  \label{tab:tornado_reg} 
\begin{tabular}{@{\extracolsep{5pt}}lcc} 
\\[-1.8ex]\hline 
\hline \\[-1.8ex] 
 & \multicolumn{2}{c}{\textit{Dependent variable:}} \\ 
\cline{2-3} 
\\[-1.8ex] & \multicolumn{2}{c}{Fuerza del Tornado (F)} \\ 
\\[-1.8ex] & (1) & (2) \\ 
\hline \\[-1.8ex] 
 \textit{Velocidad del Viento} & 0.025$^{***}$ & 0.020$^{***}$ \\ 
  & (0.0016) & (0.0025) \\ 
  & & \\ 
 \textit{Precipitaciones} &  & 0.030$^{***}$ \\ 
  &  & (0.0050) \\ 
  & & \\ 
 \textit{Temperatura} &  & 0.050 \\ 
  &  & (0.0300) \\ 
  & & \\ 
 Intercepto & -2.000$^{***}$ & -1.500$^{***}$ \\ 
  & (0.289) & (0.390) \\ 
  & & \\ 
\hline \\[-1.8ex] 
Observaciones & 182 & 182 \\ 
R$^{2}$ & 0.45 & 0.62 \\ 
Adjusted R$^{2}$ & 0.44 & 0.61 \\ 
Error estándar residual & 0.45 (df = 180) & 0.35 (df = 178) \\ 
F Statistic & 148.0$^{***}$ (df = 1; 180) & 96.3$^{***}$ (df = 3; 178) \\ 
\hline 
\hline \\[-1.8ex] 
\textit{Nota:}  & \multicolumn{2}{r}{$^{*}$p$<$0.1; $^{**}$p$<$0.05; $^{***}$p$<$0.01} \\
\textit{}  & \multicolumn{2}{r}{Valores en paréntesis corresponden al error estándar}
\end{tabular} 
\end{table}  

\begin{enumerate}
    \item \textbf{(5 puntos)} En términos de los parámetros del modelo, establezca la hipótesis nula de que, controlando por \textit{Precipitaciones} y \textit{Temperatura}, \textit{Velocidad} no tiene efecto en la \textit{Fuerza del Tornado} y luego comente la aseveración de lo anterior.
    
   {\color{blue} Para hacer esto debemos usar la información de la regresión (2), dado que queremos analizar el efecto de \textit{Velocidad} controlando por \textit{Precipitaciones} y \textit{Temperatura}. 
   
   \begin{eqnarray*}
       H_0: \gamma_1=0\\
       H_1: \gamma_1\neq 0
   \end{eqnarray*}

   Construimos el estadístico $t$ para testear nuestra hipótesis y luego lo comparamos con el valor crítico de la distribución $t_{182-3-1}$, es decir, $t_{178}$, definiendo el nivel de confianza $\alpha=0.1, 0.05, 0.01$.

    \begin{eqnarray*}
        t = \frac{0.02}{0.0025} = 8.000
    \end{eqnarray*}

    Obtenemos un valor muy alto del estadístico $t$, lo que es evidencia en contra de la hipótesis nula. Para una $t$ con 178 grados de libertad, el valor crítico con $\alpha = 0.01$ es aproximadamente $2.5760$ (normal). Por lo tanto, se rechaza la hipótesis nula al $1\%$ de confianza.
    Luego, lo que dice el enunciado es falso, y la \textit{Velocidad del viento} sí tiene efecto en la \textit{Fuerza del Tornado}.

    Como no necesitáramos reportar el valor del estadístico, podríamos haber rechazado la hipótesis de que, controlando por \textit{Precipitaciones} y \textit{Temperatura}, \textit{Velocidad del viento} no tiene efecto sobre la fuerza del tornado, al ver que en la regresión (2) el estimador de  \textit{Velocidad} tiene (***), por lo tanto, su p-valor es menor a 1\%, es decir, esa variable es estadísticamente significativa al $1\%$.}

    \item \textbf{(5 puntos)} Interprete los coeficientes en la regresión (2). En particular \textit{Velocidad del Viento}, ¿es mayor o menor que el coeficiente correspondiente en la regresión simple (1)? Explique.\\
{\color{blue}
    \begin{itemize}
    \item \textit{Velocidad del Viento}: El coeficiente estimado es $0.020$, lo que indica que, manteniendo constantes las otras variables (\textit{Precipitaciones} y \textit{Temperatura}), un incremento de 1 km/h en la velocidad del viento se asocia con un aumento promedio de $0.020$ unidades en la fuerza del tornado (medida en escala Fujita). Significativo al 1\%.
    
    \item \textit{Precipitaciones}: El coeficiente estimado es $0.030$, lo que sugiere que, manteniendo constantes las demás variables, un incremento de 1 mm en las precipitaciones se asocia con un aumento promedio de $0.030$ unidades en la fuerza del tornado. Significativo al 1\%.
    
    \item \textit{Temperatura}: El coeficiente estimado es $0.050$, indicando que un incremento de 1°C en la temperatura durante el tornado está asociado con un aumento promedio de $0.050$ unidades en la fuerza del tornado, manteniendo constantes las otras variables. No significativo.
    
    \item \textit{Intercepto}: La constante es $-1.500$, lo que representa la fuerza esperada del tornado cuando las tres variables explicativas son iguales a cero (aunque este escenario puede carecer de interpretación práctica real).
\end{itemize}

    
    \textbf{Comparación de \textit{Velocidad del Viento} entre las regresiones (1) y (2)}:

    En la regresión (1) (modelo simple), el coeficiente estimado para \textit{Velocidad del Viento} es $0.025$. Al comparar con el modelo (2) (regresión múltiple), donde el coeficiente es $0.020$, observamos que el coeficiente disminuye al incluir las variables \textit{Precipitaciones} y \textit{Temperatura}.

    Esto se explica por el hecho de que, en la regresión simple (1), el efecto de \textit{Velocidad del Viento} incluye tanto su efecto directo sobre la fuerza del tornado como efectos indirectos relacionados con las variables omitidas. En el modelo múltiple (2), esos efectos indirectos son controlados, y el coeficiente representa el efecto directo puro de \textit{Velocidad del Viento} sobre la fuerza del tornado. 

    En resumen, el coeficiente en (2) es \textbf{menor} que en (1) porque parte del efecto que en (1) se atribuía a \textit{Velocidad del Viento} en realidad está compartido con \textit{Precipitaciones} y \textit{Temperatura}.}

    \item \textbf{(5 puntos)} Testee la significancia global del modelo al 1\% y explique.\footnote{Hint: considere que el valor crítico de F al 1\% con grados de libertad (3,178) es aproximadamente 3.87}\\
{\color{blue}
    La prueba de significancia global evalúa si el modelo en su conjunto es estadísticamente significativo, es decir, si al menos una de las variables explicativas tiene un efecto diferente de cero sobre la variable dependiente.  

    \begin{itemize}
        \item La hipótesis nula (\( H_0 \)) plantea que todos los coeficientes son simultáneamente cero:
        \[
        H_0: \beta_1 = \beta_2 = \beta_3 = 0
        \]
        \item La hipótesis alternativa (\( H_1 \)) es que al menos uno de los coeficientes es distinto de cero.
    \end{itemize}

    El estadístico F del modelo (2) es:
    \[
    F = 96.3, \quad df = (3, 178)
    \]
    El valor crítico de F con \(\alpha=0.01\) y \(df=(3,178)\) es aproximadamente \(3.87\).

    Dado que \( F = 96.3 \) es mucho mayor que el valor crítico, rechazamos la hipótesis nula \( H_0 \) al 1\% de significancia.

    Entonces, el modelo es estadísticamente significativo en su conjunto, lo que indica que al menos una de las variables (\textit{Velocidad del Viento}, \textit{Precipitaciones} o \textit{Temperatura}) tiene un efecto significativo sobre la fuerza del tornado.}

    \item \textbf{(5 puntos)} Para explicar el Fuerza del Tornado, ¿es mejor el modelo restringido o el no restringido? Argumente.\\
{\color{blue}
    Para determinar si el modelo no restringido (2) es significativamente mejor que el restringido (1), realizamos una prueba F:

    \begin{itemize}
        \item Hipótesis nula (\( H_0 \)): Las restricciones son válidas, es decir, los coeficientes adicionales son iguales a cero (\( \beta_2 = \beta_3 = 0 \)).
        \item Hipótesis alternativa (\( H_1 \)): Al menos uno de los coeficientes adicionales es distinto de cero.
    \end{itemize}
    
    El estadístico F se calcula como:
    \[
    F = \frac{(R^2_{UR} - R^2_{R})/j}{(1 - R^2_{UR})/(n - k - 1)}
    \]
    donde:
    \begin{itemize}
        \item \( R^2_{UR} = 0.62 \) (no restringido)
        \item \( R^2_{R} = 0.45 \) (restringido)
        \item \( j = 2 \) (número de restricciones: coeficientes adicionales en (2))
        \item \( n = 182 \) (número de observaciones)
        \item \( k = 3 \) (número de variables explicativas en el modelo no restringido)
    \end{itemize}
    
    Sustituimos los valores:
    \[
    F = \frac{(0.62 - 0.45)/2}{(1 - 0.62)/(182 - 3 - 1)} = \frac{0.17/2}{0.38/178} = \frac{0.085}{0.00213} \approx 39.91
    \]
    
    \begin{itemize}
        \item Valor crítico: Para \( j=2 \) y \( n-k-1=178 \) grados de libertad, con \(\alpha=0.01\), el valor crítico al 1\% es aproximadamente \( F_{2,178} \approx 4.72 \).
    \end{itemize}
    
    Por lo tanto, dado que \( F \approx 39.91 \gg 4.72 \), se rechaza la hipótesis nula al 1\% de significancia. Es decir, el modelo no restringido (2) es significativamente mejor que el restringido (1).
}

\end{enumerate}

\item Un grupo de analistas se encuentra realizando una investigación sobre  el mercados de propiedades de la ciudad de Boston.  
Cuentan con los siguientes datos:

\begin{itemize}
    \item Precio de las propiedades (lprice en logaritmos).
    \item La cantidad de baños de las propiedades (baths). 
    \item La cantidad de habitaciones (rooms).
        \item Antiguedad de la propiedad (age en cantidad de años y agesq los años al cuadrado). 
\end{itemize}

Sus resultados se resumen en la siguiente  tabla:

% Table created by stargazer v.5.2.3 by Marek Hlavac, Social Policy Institute. E-mail: marek.hlavac at gmail.com
% Date and time: Thu, May 29, 2025 - 09:54:11 PM
\begin{table}[!htbp] \centering 
  \caption{} 
  \label{} 
\begin{tabular}{@{\extracolsep{5pt}}lcc} 
\\[-1.8ex]\hline 
\hline \\[-1.8ex] 
 & \multicolumn{2}{c}{\textit{Dependent variable:}} \\ 
\cline{2-3} 
\\[-1.8ex] & \multicolumn{2}{c}{lprice} \\ 
\\[-1.8ex] & (1) & (2)\\ 
\hline \\[-1.8ex] 
 larea & 0.437$^{***}$ & 0.545$^{***}$ \\ 
  & (0.073) & (0.069) \\ 
  & & \\ 
 baths & 0.224$^{***}$ & 0.070$^{*}$ \\ 
  & (0.034) & (0.038) \\ 
  & & \\ 
 rooms & 0.031 & 0.041$^{*}$ \\ 
  & (0.024) & (0.023) \\ 
  & & \\ 
 age &  & $-$0.009$^{***}$ \\ 
  &  & (0.002) \\ 
  & & \\ 
 agesq &  & 0.00003$^{***}$ \\ 
  &  & (0.00001) \\ 
  & & \\ 
 Constant & 7.328$^{***}$ & 6.914$^{***}$ \\ 
  & (0.490) & (0.461) \\ 
  & & \\ 
\hline \\[-1.8ex] 
Observations & 321 & 321 \\ 
R$^{2}$ &  & 0.597 \\ 
Adjusted R$^{2}$ &  & 0.590 \\ 
Residual Std. Error & 0.305 (df = 317) & 0.280 (df = 315) \\ 
F Statistic & 114.881$^{***}$ (df = 3; 317) & 93.213$^{***}$ (df = 5; 315) \\ 
\hline 
\hline \\[-1.8ex] 
\textit{Note:}  & \multicolumn{2}{r}{$^{*}$p$<$0.1; $^{**}$p$<$0.05; $^{***}$p$<$0.01} \\ 
\end{tabular} 
\end{table} 

\begin{enumerate}
    \item (6 puntos) En base a la información de la tabla determine el $R^2$ de la primera regresión.

    {\color{blue} El estadístico F de la primera regresión es \( F = 114{,}881 \), con \( df_1 = 3 \) (número de regresores) y \( df_2 = 317 \) (grados de libertad residuales). Usamos la fórmula del estadístico F para recuperar el \( R^2 \):

    \[
    F = \frac{R^2 / k}{(1 - R^2) / (n - k - 1)}
    \]

    Donde:
    \begin{itemize}
        \item \( k = 3 \): número de regresores en el modelo (sin contar la constante)
        \item \( n = 321 \): número de observaciones
        \item \( F = 114{,}881 \)
    \end{itemize}

    Sustituyendo en la fórmula:

    \[
    114{,}881 = \frac{R^2 / 3}{(1 - R^2) / 317}
    \Rightarrow 114{,}881 = \frac{317 R^2}{3 (1 - R^2)}
    \]

    Multiplicando cruzado y resolviendo:

    \[
    114{,}881 \cdot 3 (1 - R^2) = 317 R^2
    \Rightarrow 344{,}643 (1 - R^2) = 317 R^2
    \Rightarrow 344{,}643 - 344{,}643 R^2 = 317 R^2
   \]
    \[\Rightarrow 344{,}643 = (317 + 344{,}643) R^2 \quad R^2 = \frac{344{,}643}{661{,}643} \approx 0{,}521
    \]

    Por lo tanto, el \( R^2 \) del primer modelo es aproximadamente \textbf{0{,}521}.

También se pueden encontrar la suma de cuadrados totales del $R^2$ de la segunda regresión.}
    \item (6 puntos) Según la segunda regresión, si el área aumenta en 1\%, cuál es el cambio esperado en el precio de la propiedad? ¿Es esta variable explicativa estadísticamente significativa? Explique.  

    {\color{blue}  En la segunda regresión, el coeficiente de \texttt{larea} es 0{,}545. Dado que \texttt{lprice} es el logaritmo del precio y \texttt{larea} es el logaritmo del área, el coeficiente puede interpretarse como una elasticidad. Por lo tanto, un aumento del 1\% en el área se asocia con un incremento esperado de aproximadamente 0{,}545\% en el precio de la propiedad.

    Esta variable es estadísticamente significativa al 1\%, como lo indica el triple asterisco (***). Esto implica que hay suficiente evidencia estadística para afirmar que el área tiene un efecto significativo sobre el precio de la propiedad.}
    \item (8 puntos) Describa la relación estimada entre el precio de las propiedades y la antig\"uedad de ésta. Cree usted que es apropiada la inclusión de la antig\"uedad en el modelo. Justifique su respuesta.

    {\color{blue} La relación estimada entre el precio y la antigüedad de la propiedad es no lineal, ya que se incluye tanto la variable \texttt{age} como su cuadrado \texttt{agesq}. El coeficiente de \texttt{age} es negativo (\(-0{,}009\)) y el de \texttt{agesq} es positivo (\(0{,}00003\)), ambos estadísticamente significativos al 1\%.

    Esta especificación sugiere que el precio disminuye con la antigüedad a tasas decrecientes: al inicio, propiedades más antiguas tienen menor precio, pero ese efecto negativo se reduce a medida que la propiedad envejece más. Este comportamiento puede reflejar que, si bien las propiedades pierden valor con el tiempo, aquellas muy antiguas podrían tener características que les otorguen valor (por ejemplo, ubicación o valor patrimonial).

El efecto marginal de la antig\"uedad esta determinado por:

${\partial lprice \over \partial age}= −0.009+2ages 0.00003$

    \textbf{La inclusión de la antigüedad y su cuadrado es apropiada},  ambos coeficientes son significativos, lo que apoya su inclusión desde un punto de vista estadístico.

    También se puede realizar un test F que compare modelo restringido y no restringido.
    }
\end{enumerate}


\end{enumerate}
    \end{enumerate}








\newpage


\begin{center}
    {\Large Formulario  }

 
\end{center}

\begin{itemize}

% \item  
% {\color{red}
% $$\hat{\beta}_0 = \bar{y} - \hat{\beta}_1 \bar{x} $$
% $$\hat{\beta}_1 = \frac{\sum_{i=1}^{n}(x_i - \bar{x})(y_i-\bar{y})}{\sum_{i=1}^{n}(x_i - \bar{x})^2}$$}

\item  {\color{red} \[ \hat{\beta} = (X^{\top}X)^{-1}X^{\top}y \]}

\item 
	{\color{red}	\[ SST = \sum_{i=1}^{n} (y_i - \bar{y})^2 \quad SSE = \sum_{i=1}^{n} (\hat{y}_i - \bar{y})^2  \quad SSR = \sum_{i=1}^{n} \hat{u}_i^2\]}
		

    \item {\color{red} \[ R^2  = 1 - \frac{SSR}{SST} \quad \quad \bar{R}^2 =  1 - \frac{(n-1)SSR}{(n-k-1)SST} \]}

 \item {\color{red} \[ t_{\hat{\beta_j}}=\frac{\hat{\beta}_j - \beta_j}{\text{s.e.}(\hat{\beta}_j)} \sim t_{n-k-1} \]}

 \item {\color{red}\[
    F = \frac{SSR_{r} - SSR_{ur}}{SSR_{ur}} \cdot \frac{n - k - 1}{j} \sim F(j, n - k - 1)
    \]}

% \item {\color{red}Considere que el valor crítico a un nivel de significancia del 5\% es 1.645}
% \item {\color{red}Considere los siguientes valores críticos para t: $\alpha = 0.01 \implies t \approx 2.62$; $\alpha = 0.05 \implies t \approx 1,97$;$\alpha = 0.01 \implies t \approx 2.62$; $\alpha = 0.1 \implies t \approx 1,645$}
\item {\color{red}Considere que el valor crítico de F al 1\% con grados de libertad (2,178) es aproximadamente 4.72}

\end{itemize}

\begin{figure}[H]
    \centering
    \includegraphics[width=0.85\linewidth]{Ejercicios//ejercicios_inferencia/cert2_F2_1184.png}
\end{figure}

\begin{figure}[H]
    \centering
    \includegraphics[width=0.85\linewidth]{certamenes/normal.png}
\end{figure}


\begin{figure}
    \centering
    \includegraphics[width=0.85\linewidth]{certamenes/normal_1cola.png}
\end{figure}
\end{document}